% Options for packages loaded elsewhere
\PassOptionsToPackage{unicode}{hyperref}
\PassOptionsToPackage{hyphens}{url}
\documentclass[
]{article}
\usepackage{xcolor}
\usepackage[margin=1in]{geometry}
\usepackage{amsmath,amssymb}
\setcounter{secnumdepth}{-\maxdimen} % remove section numbering
\usepackage{iftex}
\ifPDFTeX
  \usepackage[T1]{fontenc}
  \usepackage[utf8]{inputenc}
  \usepackage{textcomp} % provide euro and other symbols
\else % if luatex or xetex
  \usepackage{unicode-math} % this also loads fontspec
  \defaultfontfeatures{Scale=MatchLowercase}
  \defaultfontfeatures[\rmfamily]{Ligatures=TeX,Scale=1}
\fi
\usepackage{lmodern}
\ifPDFTeX\else
  % xetex/luatex font selection
\fi
% Use upquote if available, for straight quotes in verbatim environments
\IfFileExists{upquote.sty}{\usepackage{upquote}}{}
\IfFileExists{microtype.sty}{% use microtype if available
  \usepackage[]{microtype}
  \UseMicrotypeSet[protrusion]{basicmath} % disable protrusion for tt fonts
}{}
\makeatletter
\@ifundefined{KOMAClassName}{% if non-KOMA class
  \IfFileExists{parskip.sty}{%
    \usepackage{parskip}
  }{% else
    \setlength{\parindent}{0pt}
    \setlength{\parskip}{6pt plus 2pt minus 1pt}}
}{% if KOMA class
  \KOMAoptions{parskip=half}}
\makeatother
\usepackage{color}
\usepackage{fancyvrb}
\newcommand{\VerbBar}{|}
\newcommand{\VERB}{\Verb[commandchars=\\\{\}]}
\DefineVerbatimEnvironment{Highlighting}{Verbatim}{commandchars=\\\{\}}
% Add ',fontsize=\small' for more characters per line
\usepackage{framed}
\definecolor{shadecolor}{RGB}{248,248,248}
\newenvironment{Shaded}{\begin{snugshade}}{\end{snugshade}}
\newcommand{\AlertTok}[1]{\textcolor[rgb]{0.94,0.16,0.16}{#1}}
\newcommand{\AnnotationTok}[1]{\textcolor[rgb]{0.56,0.35,0.01}{\textbf{\textit{#1}}}}
\newcommand{\AttributeTok}[1]{\textcolor[rgb]{0.13,0.29,0.53}{#1}}
\newcommand{\BaseNTok}[1]{\textcolor[rgb]{0.00,0.00,0.81}{#1}}
\newcommand{\BuiltInTok}[1]{#1}
\newcommand{\CharTok}[1]{\textcolor[rgb]{0.31,0.60,0.02}{#1}}
\newcommand{\CommentTok}[1]{\textcolor[rgb]{0.56,0.35,0.01}{\textit{#1}}}
\newcommand{\CommentVarTok}[1]{\textcolor[rgb]{0.56,0.35,0.01}{\textbf{\textit{#1}}}}
\newcommand{\ConstantTok}[1]{\textcolor[rgb]{0.56,0.35,0.01}{#1}}
\newcommand{\ControlFlowTok}[1]{\textcolor[rgb]{0.13,0.29,0.53}{\textbf{#1}}}
\newcommand{\DataTypeTok}[1]{\textcolor[rgb]{0.13,0.29,0.53}{#1}}
\newcommand{\DecValTok}[1]{\textcolor[rgb]{0.00,0.00,0.81}{#1}}
\newcommand{\DocumentationTok}[1]{\textcolor[rgb]{0.56,0.35,0.01}{\textbf{\textit{#1}}}}
\newcommand{\ErrorTok}[1]{\textcolor[rgb]{0.64,0.00,0.00}{\textbf{#1}}}
\newcommand{\ExtensionTok}[1]{#1}
\newcommand{\FloatTok}[1]{\textcolor[rgb]{0.00,0.00,0.81}{#1}}
\newcommand{\FunctionTok}[1]{\textcolor[rgb]{0.13,0.29,0.53}{\textbf{#1}}}
\newcommand{\ImportTok}[1]{#1}
\newcommand{\InformationTok}[1]{\textcolor[rgb]{0.56,0.35,0.01}{\textbf{\textit{#1}}}}
\newcommand{\KeywordTok}[1]{\textcolor[rgb]{0.13,0.29,0.53}{\textbf{#1}}}
\newcommand{\NormalTok}[1]{#1}
\newcommand{\OperatorTok}[1]{\textcolor[rgb]{0.81,0.36,0.00}{\textbf{#1}}}
\newcommand{\OtherTok}[1]{\textcolor[rgb]{0.56,0.35,0.01}{#1}}
\newcommand{\PreprocessorTok}[1]{\textcolor[rgb]{0.56,0.35,0.01}{\textit{#1}}}
\newcommand{\RegionMarkerTok}[1]{#1}
\newcommand{\SpecialCharTok}[1]{\textcolor[rgb]{0.81,0.36,0.00}{\textbf{#1}}}
\newcommand{\SpecialStringTok}[1]{\textcolor[rgb]{0.31,0.60,0.02}{#1}}
\newcommand{\StringTok}[1]{\textcolor[rgb]{0.31,0.60,0.02}{#1}}
\newcommand{\VariableTok}[1]{\textcolor[rgb]{0.00,0.00,0.00}{#1}}
\newcommand{\VerbatimStringTok}[1]{\textcolor[rgb]{0.31,0.60,0.02}{#1}}
\newcommand{\WarningTok}[1]{\textcolor[rgb]{0.56,0.35,0.01}{\textbf{\textit{#1}}}}
\usepackage{graphicx}
\makeatletter
\newsavebox\pandoc@box
\newcommand*\pandocbounded[1]{% scales image to fit in text height/width
  \sbox\pandoc@box{#1}%
  \Gscale@div\@tempa{\textheight}{\dimexpr\ht\pandoc@box+\dp\pandoc@box\relax}%
  \Gscale@div\@tempb{\linewidth}{\wd\pandoc@box}%
  \ifdim\@tempb\p@<\@tempa\p@\let\@tempa\@tempb\fi% select the smaller of both
  \ifdim\@tempa\p@<\p@\scalebox{\@tempa}{\usebox\pandoc@box}%
  \else\usebox{\pandoc@box}%
  \fi%
}
% Set default figure placement to htbp
\def\fps@figure{htbp}
\makeatother
\setlength{\emergencystretch}{3em} % prevent overfull lines
\providecommand{\tightlist}{%
  \setlength{\itemsep}{0pt}\setlength{\parskip}{0pt}}
\usepackage{bookmark}
\IfFileExists{xurl.sty}{\usepackage{xurl}}{} % add URL line breaks if available
\urlstyle{same}
\hypersetup{
  pdftitle={Culminating Activity},
  pdfauthor={Aloba, Loyan Edrian \textbar{} Dalmonte, Franklin \textbar{} Delrio Martin Nicholas \textbar{} Dorin, Homer Adriel \textbar{} Rivera, Nythan},
  hidelinks,
  pdfcreator={LaTeX via pandoc}}

\title{Culminating Activity}
\author{Aloba, Loyan Edrian \textbar{} Dalmonte, Franklin \textbar{}
Delrio Martin Nicholas \textbar{} Dorin, Homer Adriel \textbar{} Rivera,
Nythan}
\date{May 08, 2026}

\begin{document}
\maketitle

\section{Abstract}\label{abstract}

Provide a concise summary of the study including purpose, methods,
results, and conclusions (150--250 words).

\section{Chapter 1: Introduction}\label{chapter-1-introduction}

\subsection{Background of the Study}\label{background-of-the-study}

The integration of Generative AI (GenAI) into software engineering has
been aggressively marketed as an unprecedented paradigm shift in
productivity. Industry narratives frequently promise that AI coding
assistants will streamline workflows, accelerate output, and
consequently elevate both financial compensation and developer
well-being. However, emerging empirical research reveals a significant
productivity trade-off inherent in GenAI adoption. Utilizing the SPACE
framework for developer productivity, Afroz et al.~(2025) found that
while developers using AI tools experience localized spikes in activity
and perceive themselves as operating faster, they do not necessarily
produce better software or experience increased fulfillment. In fact,
rigorous randomized controlled trials indicate that developers using
early-2025 AI tools took 19\% longer to complete complex tasks, despite
subjectively believing they were working 20\% faster---a cognitive
disconnect establishing an industry-wide ``productivity placebo''
(Becker et al., 2025). This paradox is driven by severe cognitive and
operational friction. As AI accelerates initial code generation, it
shifts systemic bottlenecks downstream, increasing pull request review
times by 91\% and inflating the volume of bugs per developer (Faros AI,
2025). Developers are forced to spend more time debugging hallucinated
code and navigating complex, AI-generated architectural debt, which has
been shown to actively damage macroscopic software delivery performance
and system stability (DORA Research Program, 2024). Consequently, this
study investigates whether the financial rewards of AI adoption
genuinely translate into higher job satisfaction, or if the persistent
friction of modern coding creates a ``GenAI Productivity-Satisfaction
Tradeoff.'' In this paradigm, developers may be caught in a cycle where
hyper-compensation---driven by a reported 56\% AI skill premium
(PricewaterhouseCoopers, 2025)---acts as a retention mechanism despite
high stress, a concept previously identified in occupational psychology
as job embeddedness or ``golden handcuffs'' (Mitchell et al., 2001). As
AI increasingly shifts the developer's role from writing original code
to reviewing and debugging AI output, the industry must question whether
these tools actually improve the human experience of software
engineering. Utilizing secondary data from the Stack Overflow 2024
Developer Survey (Stack Overflow, 2024), this study builds a predictive
model to determine if AI tools genuinely improve job satisfaction, or if
developers are simply compensated more to endure increased friction.

\subsection{Statement of the Problem}\label{statement-of-the-problem}

This study investigates the systemic impact of AI adoption on software
developers. With Job Satisfaction as the primary target variable, the
study asks: To what extent does the use of AI development tools predict
a developer's overall job satisfaction? How does workflow friction (time
spent searching for answers and debugging) interact with AI usage to
influence this satisfaction? Does a higher salary offset the negative
impact of workflow friction when predicting overall job satisfaction?

\subsection{Objectives of the Study}\label{objectives-of-the-study}

\subsubsection{General Objective}\label{general-objective}

To investigate the ``GenAI Productivity-Satisfaction Tradeoff'' by
developing a predictive model that evaluates how AI tool adoption,
financial compensation, and workflow friction interact to determine a
software developer's overall job satisfaction.

\subsubsection{Specific Objectives}\label{specific-objectives}

To determine the extent to which AI tool adoption serves as a predictor
for both financial compensation and workflow friction (time spent
searching and debugging).

To identify the primary predictors of developer job satisfaction using a
Classification and Regression Tree (CART) machine learning model.

To evaluate the interaction between compensation and workflow friction,
specifically testing if a higher salary offsets the negative predictive
impact of AI-induced friction on overall job satisfaction.

\subsection{Hypotheses (if applicable)}\label{hypotheses-if-applicable}

H0 (Null Hypothesis): AI usage, financial compensation, and workflow
friction do not form a significant predictive model for developer job
satisfaction, and no distinct tradeoff exists between these variables.

H1 (Alternative Hypothesis): The optimal predictive model for Job
Satisfaction will demonstrate a significant tradeoff, where the positive
predictive value of high financial compensation is canceled out by the
negative predictive impact of AI-induced workflow friction.

\subsection{Significance of the Study}\label{significance-of-the-study}

This research benefits technology companies, engineering managers, and
developers. By identifying the gap between productivity metrics such as
salary and output, and human metrics such as satisfaction, organizations
can re-evaluate how they integrate AI. It shifts the focus from simply
buying AI licenses to actively reducing the systemic friction that
causes developer burnout.

\subsection{Scope and Limitations}\label{scope-and-limitations}

The study analyzes secondary data from the Stack Overflow 2024 Developer
Survey. The scope is strictly limited to respondents who identify as
``Employed, full-time'' to maintain a stable economic baseline.

A critical limitation is that survey data regarding compensation and
satisfaction is highly susceptible to MNAR (Missing Not At Random) bias,
as extreme earners and highly burned-out developers frequently skip
these questions. To mitigate this, advanced algorithmic imputation using
MICE via Classification and Regression Trees was utilized to predict and
fill missing values rather than relying on listwise deletion.

This study does not attempt an exhaustive feature selection of the
entire Stack Overflow dataset. Instead, the predictor variables (AI
Usage, Salary, Time Spent Searching, Years of Experience, and Remote
Status) were handpicked based on intuition and prior literature
regarding developer productivity.

\section{Chapter 2: Review of Related Literature (@Loyan \&
@Nythan)}\label{chapter-2-review-of-related-literature-loyan-nythan}

\subsection{Related Studies}\label{related-studies}

Summarize previous research.

\subsection{Related Literature}\label{related-literature}

Recent work by Stephany and Teutloff (2024) demonstrates that the
economic value of a skill is strongly determined by \emph{skill
complementarity}---the number, diversity, and value of other skills with
which it can be combined. Their analysis of 962 skills reveals that AI
skills are particularly valuable, increasing worker wages by 21\% on
average, primarily because they exhibit high levels of complementarity
with diverse existing skill sets. This finding provides theoretical
grounding for understanding why financial compensation emerges as a
critical mediating variable in predicting developer satisfaction. Rather
than being arbitrary or divorced from productivity metrics, the salary
premium associated with AI adoption reflects structural economic value
accruing through skill portfolio enhancement.

\subsection{Conceptual Framework}\label{conceptual-framework}

At the core of this framework is the ultimate target variable: Job
Satisfaction.

Rather than utilizing isolated linear regressions, the framework employs
a Decision Tree architecture where predictor variables interact
organically:

Primary Predictor: AI Tool Adoption is positioned as the central
independent variable, hypothesized to simultaneously impact financial
compensation and workflow friction.

Interacting Predictors: Financial Compensation (Salary) and Workflow
Friction (Time Spent Searching) act as mediating variables within the
model. Contextual Controls: Years of Experience and Remote Status serve
as baseline variables to account for demographic and workplace
variances.

The CART algorithm evaluates how these variables interact, specifically
testing whether high compensation mathematically buffers the negative
impact of AI-induced workflow friction when predicting final job
satisfaction levels.

\begin{figure}
\includegraphics[width=1\linewidth]{Culminating-Activity_files/figure-latex/conceptual-framework-1} \caption{Conceptual Framework of the GenAI Productivity-Satisfaction Tradeoff}\label{fig:conceptual-framework}
\end{figure}

\subsection{Chapter 3: Methodology}\label{chapter-3-methodology}

\textbf{Research Design}

This study uses a quantitative, descriptive-correlational design to
examine how AI adoption, salary, and workflow friction (as measured by
Time Spent Searching \& Debugging in minutes) relate to developer job
satisfaction.

\textbf{Data Collection} The study uses secondary data from the Stack
Overflow 2024 Developer Survey. The raw CSV is provided in the
repository \url{data/survey_results_public.csv}; the analysis filters
this file to full-time respondents (see \texttt{filter\_columns.r}) and
writes the reduced table to \url{outputs/data_filtered_columns.csv}.

The survey schema is included as \url{data/survey_results_schema.csv}.

Compensation and satisfaction responses are vulnerable to MNAR (Missing
Not At Random) bias, so the analysis uses MICE with Classification and
Regression Trees instead of listwise deletion.

The final modeling set was limited to Salary, Uses\_AI, Years\_Exp,
Job\_Satisfaction, Is\_Remote, and TimeSearching because these variables
align directly with compensation, AI adoption, seniority, work
arrangement, and workflow friction.

\textbf{Data Preparation and Initial Variable Selection}

The raw survey file was reduced with \url{filter_columns.r}, which
filtered for full-time employees and built the analysis table. Salary
was derived from \texttt{ConvertedCompYearly}, Years\_Exp from
\texttt{YearsCodePro}, Uses\_AI from \texttt{AISelect},
Job\_Satisfaction from \texttt{JobSat}, and Is\_Remote from
\texttt{RemoteWork}. The output was saved to
\url{outputs/data_filtered_columns.csv}.

These columns were kept because they capture the financial, behavioral,
and contextual factors needed for this analysis' statement of the
problem and objectives. Salary was prioritized as a key predictor
because prior research on skill complementarity (Stephany \& Teutloff,
2024) establishes that compensation reflects structural economic value
rather than arbitrary scarcity premiums, making it a theoretically
grounded mediating variable between AI adoption and job satisfaction.

\textbf{Initial Data Checks and EDA}

The dataset was checked with \url{check_data.r} to review missingness
patterns, unique-value counts, and basic distributions in both the raw
survey data and the filtered modeling table.

Exploratory visualization was handled with \url{eda.rmd} to inspect the
early relationships among AI adoption, salary, job satisfaction,
experience, and time spent searching/debugging. This step served as a
plausibility check before modeling.

\textbf{Imputation and Modeling Workflow}

The imputation step was handled in \url{main.r}, which creates or reuses
\url{outputs/data_imputed_primary.csv}.

An initial exploratory run used \texttt{set.seed(42)} (a commonly used
placeholder seed; see Appendix C). Because a single RNG draw can create
fragile results, the pipeline includes an explicit seed-sensitivity
check: imputations were re-run with seeds = c(42, 100, 200, 300, 400)
using \texttt{cart\_robustness\_check.r}, and the resulting trees were
compared by their cross-validated error (\texttt{xerror}). This
procedure ensures we report a defensible baseline while keeping the
imputation uncertainty transparent.

After the sensitivity check, subsequent script runs adopted
\texttt{set.seed(200)} for the corrected baseline because it produced
the lowest cross-validated error across the tested seeds (see Chapter 4
and \url{outputs/cart_mice_seed_averages.csv}).

\textbf{Hyperparameter Shuffling}

The CART stage then uses hyperparameter shuffling to compare
\texttt{cp}, \texttt{minsplit}, and \texttt{maxdepth}. In this context,
\texttt{cp} is the pruning strength and \texttt{xerror} is the
cross-validated error reported by \texttt{rpart}. Higher \texttt{cp}
values prune more aggressively and can remove all splits. The
best-performing combination was \texttt{cp\ =\ 0.001},
\texttt{minsplit\ =\ 20}, and \texttt{maxdepth\ =\ 5} (see
hyperparameter outputs), and under the corrected baseline
\texttt{TimeSearching\_Min} was the strongest predictor, followed by
\texttt{Years\_Exp} and \texttt{Salary}.

\textbf{Robustness Checks}

The robustness work was separated into two parts. First, the
hyperparameter shuffling held the imputed data fixed while varying model
settings, which isolates model-choice effects. Second, the
seed-sensitivity check changed the imputation seed while holding the
CART settings fixed, which isolates data-uncertainty effects.

The seed-sensitivity workflow is documented in
\url{cart_robustness_check.r}, and its outputs are saved to
\url{outputs/cart_mice_seed_runs.csv} and
\url{outputs/cart_mice_seed_averages.csv}. The CP profile is saved to
\url{outputs/cart_cp_profile.csv} and provides an additional check on
pruning.

\subsection{Chapter 4: Results and
Discussion}\label{chapter-4-results-and-discussion}

\textbf{Descriptive Statistics}

\begin{Shaded}
\begin{Highlighting}[]
\FunctionTok{library}\NormalTok{(dplyr)}
\end{Highlighting}
\end{Shaded}

\begin{verbatim}
## 
## Attaching package: 'dplyr'
\end{verbatim}

\begin{verbatim}
## The following objects are masked from 'package:stats':
## 
##     filter, lag
\end{verbatim}

\begin{verbatim}
## The following objects are masked from 'package:base':
## 
##     intersect, setdiff, setequal, union
\end{verbatim}

\begin{Shaded}
\begin{Highlighting}[]
\FunctionTok{library}\NormalTok{(readr)}

\NormalTok{df }\OtherTok{\textless{}{-}} \FunctionTok{read\_csv}\NormalTok{(}\StringTok{"outputs/data\_imputed\_primary.csv"}\NormalTok{, }\AttributeTok{show\_col\_types =} \ConstantTok{FALSE}\NormalTok{)}

\NormalTok{summary\_df }\OtherTok{\textless{}{-}}\NormalTok{ df }\SpecialCharTok{\%\textgreater{}\%}
  \FunctionTok{summarise}\NormalTok{(}
    \AttributeTok{n =} \FunctionTok{n}\NormalTok{(),}
    \AttributeTok{Job\_Satisfaction\_mean =} \FunctionTok{mean}\NormalTok{(Job\_Satisfaction, }\AttributeTok{na.rm =} \ConstantTok{TRUE}\NormalTok{),}
    \AttributeTok{Job\_Satisfaction\_sd =} \FunctionTok{sd}\NormalTok{(Job\_Satisfaction, }\AttributeTok{na.rm =} \ConstantTok{TRUE}\NormalTok{),}
    \AttributeTok{Salary\_mean =} \FunctionTok{mean}\NormalTok{(Salary, }\AttributeTok{na.rm =} \ConstantTok{TRUE}\NormalTok{),}
    \AttributeTok{Salary\_sd =} \FunctionTok{sd}\NormalTok{(Salary, }\AttributeTok{na.rm =} \ConstantTok{TRUE}\NormalTok{),}
    \AttributeTok{TimeSearching\_Min\_mean =} \FunctionTok{mean}\NormalTok{(TimeSearching\_Min, }\AttributeTok{na.rm =} \ConstantTok{TRUE}\NormalTok{),}
    \AttributeTok{TimeSearching\_Min\_sd =} \FunctionTok{sd}\NormalTok{(TimeSearching\_Min, }\AttributeTok{na.rm =} \ConstantTok{TRUE}\NormalTok{),}
    \AttributeTok{Years\_Exp\_mean =} \FunctionTok{mean}\NormalTok{(Years\_Exp, }\AttributeTok{na.rm =} \ConstantTok{TRUE}\NormalTok{),}
    \AttributeTok{Years\_Exp\_sd =} \FunctionTok{sd}\NormalTok{(Years\_Exp, }\AttributeTok{na.rm =} \ConstantTok{TRUE}\NormalTok{),}
    \AttributeTok{Uses\_AI\_rate =} \FunctionTok{mean}\NormalTok{(}\FunctionTok{as.numeric}\NormalTok{(}\FunctionTok{as.character}\NormalTok{(Uses\_AI)) }\SpecialCharTok{==} \DecValTok{1}\NormalTok{, }\AttributeTok{na.rm =} \ConstantTok{TRUE}\NormalTok{),}
    \AttributeTok{Is\_Remote\_rate =} \FunctionTok{mean}\NormalTok{(}\FunctionTok{as.numeric}\NormalTok{(}\FunctionTok{as.character}\NormalTok{(Is\_Remote)) }\SpecialCharTok{==} \DecValTok{1}\NormalTok{, }\AttributeTok{na.rm =} \ConstantTok{TRUE}\NormalTok{)}
\NormalTok{  )}

\NormalTok{summary\_df}
\end{Highlighting}
\end{Shaded}

\begin{verbatim}
## # A tibble: 1 x 11
##       n Job_Satisfaction_mean Job_Satisfaction_sd Salary_mean Salary_sd
##   <int>                 <dbl>               <dbl>       <dbl>     <dbl>
## 1 39041                  6.93                2.08      88450.   122471.
## # i 6 more variables: TimeSearching_Min_mean <dbl>, TimeSearching_Min_sd <dbl>,
## #   Years_Exp_mean <dbl>, Years_Exp_sd <dbl>, Uses_AI_rate <dbl>,
## #   Is_Remote_rate <dbl>
\end{verbatim}

Visualization (@Martin)

\begin{Shaded}
\begin{Highlighting}[]
\CommentTok{\# Potential visualization: a compact bar chart or table snapshot of summary\_df.}
\end{Highlighting}
\end{Shaded}

The descriptive table above is computed from
\url{outputs/data_imputed_primary.csv}.

Sample lang to (@Martin)

The imputed analysis sample provides the descriptive base for the rest
of the chapter. Salary, job satisfaction, time spent
searching/debugging, and years of experience are the central quantities
because they capture the compensation, workload, and career-stage
dimensions of a software engineer in the industry.

\textbf{Model Complexity and Pruning Stability --- Complexity Profile
(CP)}

Before evaluating hyperparameter choices, we first establish the
fundamental complexity--accuracy tradeoff through the CP profile. This
analysis shows how the pruning threshold constrains tree size and
prediction error.

\begin{Shaded}
\begin{Highlighting}[]
\NormalTok{cp\_profile }\OtherTok{\textless{}{-}} \FunctionTok{read\_csv}\NormalTok{(}\StringTok{"outputs/cart\_cp\_profile.csv"}\NormalTok{, }\AttributeTok{show\_col\_types =} \ConstantTok{FALSE}\NormalTok{)}

\NormalTok{cp\_profile}
\end{Highlighting}
\end{Shaded}

\begin{verbatim}
## # A tibble: 6 x 5
##        CP nsplit `rel error` xerror    xstd
##     <dbl>  <dbl>       <dbl>  <dbl>   <dbl>
## 1 0.0118       0       1      1.00  0.00876
## 2 0.00553      1       0.988  0.988 0.00862
## 3 0.00267      2       0.983  0.983 0.00861
## 4 0.00141      3       0.980  0.981 0.00855
## 5 0.00116      4       0.979  0.980 0.00854
## 6 0.001        5       0.977  0.979 0.00853
\end{verbatim}

\begin{Shaded}
\begin{Highlighting}[]
\CommentTok{\# Potential visualization: line plot of xerror against CP from outputs/cart\_cp\_profile.csv.}
\end{Highlighting}
\end{Shaded}

The CP profile in \url{outputs/cart_cp_profile.csv} supports the same
conclusion but from another angle. In a decision tree algorithm, the CP
balances predictive accuracy against structural complexity, effectively
preventing the model from making unnecessary, hyper-granular splits that
lead to overfitting.

An analysis of the CP profile shows a clear stabilization point. The
table in \url{outputs/cart_cp_profile.csv} reports xerror values of:
1.00002 at CP = 0.0118428 (nsplit = 0), 0.98821 at CP = 0.0055319
(nsplit = 1), 0.98347 at CP = 0.0026676 (nsplit = 2), 0.98070 at CP =
0.00140918 (nsplit = 3), 0.98004 at CP = 0.00115968 (nsplit = 4), and
0.97894 at CP = 0.001 (nsplit = 5). The largest reductions in xerror
occur with the earliest splits; after CP = 0.001 the marginal
improvements are very small.

As seen above the cross-validated error across different pruning
thresholds reveals a clear stabilization point at CP = 0.001 As the tree
begins to grow, the initial broad splits, which represent the
macro-structural drivers of Experience and Salary, produce the most
significant and reliable reductions in prediction error. However, as the
complexity constraint is loosened to allow for deeper branches, the
cross-validated error flattens out and changes only within a very narrow
range.

This behavior makes the chosen tree complexity highly defensible: the
model attains most of its predictive benefit with a small number of
splits, capturing the principal signals (years of experience and salary)
without overfitting to minor fluctuations in workflow friction (time
spent searching/debugging) and/or the adoption of AI tools.

\textbf{Hyperparameter Shuffling}

With the CP stabilization point at 0.001 identified, we now examine
which hyperparameter combination achieves the best balance of accuracy
and interpretability.

\begin{Shaded}
\begin{Highlighting}[]
\NormalTok{shuffle }\OtherTok{\textless{}{-}} \FunctionTok{read\_csv}\NormalTok{(}\StringTok{"outputs/cart\_hyperparam\_shuffle.csv"}\NormalTok{, }\AttributeTok{show\_col\_types =} \ConstantTok{FALSE}\NormalTok{)}

\NormalTok{best\_shuffle }\OtherTok{\textless{}{-}}\NormalTok{ shuffle }\SpecialCharTok{\%\textgreater{}\%}
  \FunctionTok{filter}\NormalTok{(status }\SpecialCharTok{==} \StringTok{"ok"}\NormalTok{) }\SpecialCharTok{\%\textgreater{}\%}
  \FunctionTok{arrange}\NormalTok{(xerror, }\FunctionTok{desc}\NormalTok{(nsplit)) }\SpecialCharTok{\%\textgreater{}\%}
  \FunctionTok{slice}\NormalTok{(}\DecValTok{1}\NormalTok{)}

\NormalTok{best\_shuffle}
\end{Highlighting}
\end{Shaded}

\begin{verbatim}
## # A tibble: 1 x 12
##    rank status    cp minsplit maxdepth xerror nsplit Uses_AI Salary
##   <dbl> <chr>  <dbl>    <dbl>    <dbl>  <dbl>  <dbl>   <dbl>  <dbl>
## 1     5 ok     0.001       20        5  0.981    125       0   215.
## # i 3 more variables: TimeSearching_Min <dbl>, Years_Exp <dbl>, Is_Remote <dbl>
\end{verbatim}

Hyperparameter Shuffling results are taken from
\url{outputs/cart_hyperparam_shuffle.csv}. The best usable setting is cp
= 0.001, minsplit = 20, maxdepth = 5 (xerror = 0.98105, nsplit = 125).

By contrast, larger cp values (e.g., 0.008--0.02) frequently produce
nsplit = 0 --- a textbook sign of over-pruning. Over-pruning happens
when the pruning penalty (cp) requires a splitting gain larger than any
candidate split provides, so rpart collapses the tree to its root. The
practical effects are underfitting (the model cannot partition the
data), loss of variable‑importance information, and misleading
simplicity: a pruned tree with nsplit = 0 may appear stable but offers
no explanatory partitions. Use the CP profile and grid search to pick cp
near the elbow (here ≈ 0.001), where xerror is low and nsplit
\textgreater{} 0, balancing interpretability and generalization.

\begin{Shaded}
\begin{Highlighting}[]
\CommentTok{\# Potential visualization: xerror by cp, with nsplit shown as a point label or second axis.}
\end{Highlighting}
\end{Shaded}

\textbf{Set Seed Sensitivity}

Next, we verify that the main results are not artifacts of a single
imputation draw by testing sensitivity across different random seeds.

\begin{Shaded}
\begin{Highlighting}[]
\NormalTok{seed\_averages }\OtherTok{\textless{}{-}} \FunctionTok{read\_csv}\NormalTok{(}\StringTok{"outputs/cart\_mice\_seed\_averages.csv"}\NormalTok{, }\AttributeTok{show\_col\_types =} \ConstantTok{FALSE}\NormalTok{)}

\NormalTok{seed\_averages }\SpecialCharTok{\%\textgreater{}\%} \FunctionTok{arrange}\NormalTok{(xerror)}
\end{Highlighting}
\end{Shaded}

\begin{verbatim}
## # A tibble: 5 x 8
##    seed nsplit xerror Uses_AI Salary TimeSearching_Min Years_Exp Is_Remote
##   <dbl>  <dbl>  <dbl>   <dbl>  <dbl>             <dbl>     <dbl>     <dbl>
## 1   200    185  0.978   0       652.             2503.     1102.      1.23
## 2   100    192  0.979   0       504.             2445.     1111.      4.56
## 3    42    159  0.979   2.75    343.             2510.     1113.      2.18
## 4   300    166  0.979   0       502.             2487.      971.      1.76
## 5   400    186  0.979   0.632   684.             2503.      998.      3.14
\end{verbatim}

\begin{Shaded}
\begin{Highlighting}[]
\CommentTok{\# Potential visualization: variable importance by seed, using outputs/cart\_mice\_seed\_averages.csv.}
\end{Highlighting}
\end{Shaded}

The MICE averages in \url{outputs/cart_mice_seed_averages.csv} show that
the model is not driven by a single imputation draw. Across the tested
seeds, TimeSearching\_Min stays the strongest signal, while Years\_Exp
and Salary remain secondary. Uses\_AI stays weak or near zero in most
runs, which suggests that the direct effect of AI adoption is not stable
once missing data is re-generated.

In this study the lowest observed \texttt{xerror} across the tested
seeds occurred at \texttt{set.seed(200)} (xerror ≈ 0.9784), so the main
results are reported using the imputed dataset generated with
\texttt{set.seed(200)}. The full seed-run results are available in
\url{outputs/cart_mice_seed_runs.csv} and summarized in
\url{outputs/cart_mice_seed_averages.csv}.

\textbf{Testing for Robustness}

The corrected baseline does not support AI usage as a strong driver of
job satisfaction. In the main tree, TimeSearching\_Min is the dominant
split variable, with Years\_Exp and Salary following behind. This means
the model is showing that workflow friction is the strongest signal in
the fixed-imputation run.

The robustness checks explain why this matters. When the model settings
are shuffled, the same pattern stays visible: the chosen tree still
favors TimeSearching\_Min under the best-performing settings, while AI
usage remains weak and unstable. When the missing data is perturbed
through the seed sensitivity check, the same broad conclusion holds.
TimeSearching\_Min stays dominant, Years\_Exp and Salary remain
secondary, and Uses\_AI contributes little or nothing in most runs.

The real drivers of satisfaction appear to be workflow friction, years
of experience, and salary. This is closer to Herzberg's (1959) view that
satisfaction rests on basic working conditions and compensation, not on
a single new tool. AI may still matter in limited cases, but it is not
the main driver of the tree. In the corrected baseline, the model points
first to the time spent searching and debugging, then to experience and
pay.

This means GenAI is not a universal fix for low morale. It behaves more
like a conditional feature than a core predictor. The 56\% wage premium
for AI skills (PricewaterhouseCoopers, 2025) can still work like
``golden handcuffs'' (Mitchell et al., 2001), but the tree does not show
AI usage itself as the central force shaping satisfaction.

Interpretation

Taken together, the evidence suggests that the main tree is stable in
its structure at the chosen hyperparameters, but the importance ranking
is sensitive to imputation. In the corrected baseline,
TimeSearching\_Min is the strongest predictor, with Years\_Exp and
Salary trailing behind. When the missing data is perturbed through
different imputations, workflow friction remains the top signal, while
Uses\_AI stays weak and unstable.

This shows a weak and non-robust direct link between AI use and job
satisfaction, while workflow friction, years of experience, and salary
remain the strongest and most consistent signals.

\subsection{Chapter 5: Conclusion and
Recommendations}\label{chapter-5-conclusion-and-recommendations}

Conclusion

This research investigated the systemic impact of AI adoption on
software developers, specifically examining whether the psychological
burden of AI-induced workflow friction negates the positive impacts of
financial compensation. Addressing the core objectives established at
the outset of the study, the findings reveal that the direct link
between AI tool adoption and overall job satisfaction is surprisingly
weak. While the initial problem statement questioned how workflow
friction interacts with AI usage, the data demonstrates that at a broad
career level, developers look past this friction. This behavior is
driven by a ``productivity placebo,'' where the perception of speed
outweighs the actual workload (Becker et al., 2025). However, at a
granular operational level, the daily time spent debugging remains a
volatile source of frustration, confirming the warnings from Afroz et
al.~(2025) regarding the hidden effort of reviewing AI outputs.
Regarding the initial inquiry of whether a higher salary offsets this
friction, the findings point to a structural reality rather than a
direct mathematical tradeoff. High compensation acts as ``golden
handcuffs'' (Mitchell et al., 2001), successfully keeping developers in
their roles. This wage premium reflects the structural economic value of
AI skills through complementarity with diverse existing competencies
(Stephany \& Teutloff, 2024), making the financial incentive rational
from both employer and employee perspectives. The daily annoyance of
technical debt is endured because the financial compensation makes it
rational to stay, rather than the friction being directly ``canceled
out'' by the pay. These findings align directly with the observations by
Afroz et al.~(2025): faster task completion through GenAI does not
inherently represent true progress. Organizations should avoid the
assumption that integrating AI coding assistants will automatically
yield higher productivity or a happier workforce.

The data clearly indicates that GenAI can create ``spurious
productivity,'' where surface-level acceleration hides the deeper,
manual effort required to review and validate AI outputs (Afroz et al.,
2025).

Consequently, the alternative hypothesis (H1), which states that there
is a significant tradeoff in which high financial compensation is
canceled out by AI-induced friction, is definitively rejected (See
Chapter 4 for Set Seed Sensitivity and Testing for Robustness
subsections, where Uses\_AI remains weak or near zero in most runs). The
anticipated tradeoff was not reflected in the analysis of the dataset,
because the predictive power of AI usage is not reflected in the
different robustness checks. Overall job satisfaction is driven more by
foundational factors like years of experience and competitive pay, while
workflow friction operates independently as an expected daily challenge.
AI adoption remains a highly conditional modifier, not a universal
driver of a developer's professional well-being.

\textbf{Recommendations}

The following recommendations are intended for future researchers who
want to extend or replicate this study's method. They are grounded in
the patterns observed in Chapter 4, especially the sensitivity of the
tree to imputation and the recurring importance of workflow friction
over AI usage.

\begin{enumerate}
\def\labelenumi{\arabic{enumi}.}
\item
  Compare imputation strategies directly. MICE should remain the primary
  benchmark because it handles mixed variable types well and preserves
  uncertainty across multiple draws, but future studies can test
  KNN-based imputation as a sensitivity check. If a clustering-based
  alternative is explored, it should be used cautiously and justified
  against the data structure, since not all clustering methods handle
  mixed categorical and numeric variables cleanly. The main goal is to
  determine whether the same findings hold under a different
  missing-data treatment.
\item
  Use the imputation comparison to test robustness, not just accuracy. A
  future paper could run the same modeling pipeline under MICE, KNN, and
  one additional baseline method, then compare the downstream effects on
  tree structure, xerror, and variable importance. That would show
  whether the conclusions are stable or whether they depend too heavily
  on one imputation choice.
\item
  Extend the analysis to salary as a second outcome. A useful follow-up
  would be to build a parallel model for Salary and then compare it with
  the Job Satisfaction model. Since the results and discussion suggest
  that compensation and satisfaction are related but not identical
  outcomes, this would let future researchers check whether the same
  predictors appear in both models and whether AI tool adoption has a
  similar or different role depending on the response variable.
\item
  Make AI adoption more explicit in future models. The results section
  shows that AI usage was not the dominant driver in the final tree, so
  future researchers should still keep it in the model but consider
  stronger operationalizations, such as frequency of use, task-specific
  use, or usage intensity. This would help determine whether the null
  result comes from the construct itself or from the way it was
  measured.
\item
  Report seed and model stability in every replication. Because
  imputation and tree fitting are both stochastic, future work should
  report the seed used for the main run, the full seed-sensitivity
  table, and the selected hyperparameters. A small appendix of alternate
  runs is often enough to show that the result is not a one-seed
  artifact.
\item
  If time allows, compare CART with at least one alternative model. CART
  was appropriate here because the goal was interpretability, but future
  studies could compare it with random forest or a regularized
  regression model to check whether the same predictors appear
  consistently across different model families.
\item
  Add a cross-outcome association check when the workflow is expanded.
  The results and discussion make it reasonable to expect overlap
  between the Salary and Job Satisfaction models, so future researchers
  can formally test that relationship using a concordance check,
  correlation screen, or two-model comparison. That would be a clean way
  to see whether AI adoption, experience, and workflow friction line up
  across outcomes without forcing the idea into the current study's
  already-complete results section.
\end{enumerate}

\subsection{References}\label{references}

Afroz, S., Feng, Z., Kimura, K., Trinkenreich, B., Steinmacher, I., \&
Sarma, A. (2025). Developer productivity with GenAI. arXiv.
\url{https://arxiv.org/html/2510.24265v1}

Becker, J., Rush, N., Cunningham, T., \& Rein, D. (2025). Measuring the
impact of early-2025 AI on experienced open-source developer
productivity. arXiv. \url{https://arxiv.org/abs/2507.09089}

DORA Research Program. (2024). 2024 Accelerate state of DevOps report.
Google Cloud. \url{https://dora.dev/research/2024/dora-report/}

Faros AI. (2025). The AI productivity paradox research report.
\url{https://www.faros.ai/blog/ai-software-engineering}

Herzberg, F. (1959). The motivation to work. John Wiley \& Sons.

Mitchell, T. R., Holtom, B. C., Lee, T. W., Sablynski, C. J., \& Erez,
M. (2001). Why people stay: Using job embeddedness to predict voluntary
turnover. Academy of Management Journal, 44(6), 1102-1121.
\url{https://doi.org/10.2307/3069391}

PricewaterhouseCoopers. (2025). 2025 Global AI jobs barometer.
\url{https://www.pwc.com/gx/en/issues/artificial-intelligence/job-barometer/2025/report.pdf}

Stack Overflow. (2024). 2024 Developer Survey.
\url{https://survey.stackoverflow.co/2024/}

Stephany, F., \& Teutloff, O. (2024). What is the price of a skill? The
value of complementarity. \emph{Research Policy}, 53(1), 104927.
\url{https://doi.org/10.1016/j.respol.2023.104927}

\subsection{Appendices}\label{appendices}

Appendix A: Code (@Nythan)

The analysis code is included in the code chunks throughout this R
Markdown document.

Appendix B: Data Source Stack Overflow. (2024). 2024 Developer Survey.
Retrieved from \url{https://survey.stackoverflow.co/2024/}

Appendix C: Notes on Random Seeds

Random seeds initialize pseudorandom number generators (PRNGs) so that
stochastic processes (like MICE imputations or algorithmic subsampling)
can be reproduced exactly. A seed is not a model parameter --- it only
fixes the RNG state for reproducibility. Because different seeds can
produce slightly different imputations or model splits, reporting the
seed and testing sensitivity across several seeds is good practice for
transparency.

Why use \texttt{42}? The number 42 is a cultural reference (Douglas
Adams' The Hitchhiker's Guide to the Galaxy) and has become a de-facto,
tongue-in-cheek default in many code examples. There is no statistical
or computational reason to prefer 42 over any other integer --- it is
purely conventional.

Recommendation: report the seed used for the published baseline, include
the full seed-sensitivity results (as found in
\url{outputs/cart_mice_seed_averages.csv}), and treat the seed choice as
part of the reproducibility metadata rather than a hyperparameter that
needs theoretical justification.

\end{document}
